\chapter{Stool Color Changes}

\section*{Definition}
Stool color changes refer to variations in the color of bowel movements that can be normal or indicate underlying disease. Stool color is influenced by diet, medications, and the presence of blood or bile.

\section*{Pathophysiology}
Normal brown stool color comes from stercobilin, a breakdown product of bilirubin from red blood cell metabolism. Changes in stool color can result from: absence of bile (pale/clay-colored suggesting biliary blockage), rapid transit (green from bile not being metabolized), upper GI bleeding (black, tarry from digested blood), lower GI bleeding (bright red from fresh blood), or dietary pigments.

\section*{Causes}
\begin{itemize}
\item Black/tarry (dark, tarry stools): upper GI bleeding (ulcer, gastritis, varices), iron supplements, bismuth (Pepto-Bismol), black licorice, blueberries
\item Bright red blood (hematochezia): hemorrhoids, anal fissure, diverticular bleeding, colorectal cancer, inflammatory bowel disease
\item Pale/clay-colored: biliary blockage (gallstone, pancreatic cancer, hepatitis), barium contrast
\item Green: leafy green vegetables, green food coloring, rapid intestinal transit (diarrhea)
\item Yellow/greasy: malabsorption (celiac disease, pancreatic insufficiency, giardiasis)
\item Orange: foods high in beta-carotene (carrots, sweet potatoes), antacids, rifampin
\end{itemize}

\section*{When to See a Doctor}
\begin{itemize}
\item Black, tarry, or maroon stools (requires urgent evaluation)
\item Bright red blood on or in stool
\item Persistent pale or clay-colored stools
\item Stool color change persisting more than a few days
\item Accompanying symptoms: abdominal pain, weight loss, yellowing of the skin or eyes, fever
\item Change in bowel habits or stool caliber
\end{itemize}

\section*{Related Symptoms}
\begin{itemize}
\item Rectal bleeding
\item Gastrointestinal bleeding
\item yellowing of the skin or eyes
\item Malabsorption
\item Constipation
\item Diarrhea
\end{itemize}

\section*{Self-Care / Home Management}
\begin{itemize}
\item Review dietary intake for possible food-related causes
\item Review medications and supplements
\item If black stools from iron or bismuth, discontinue and observe
\item Increase fiber and water intake
\item Keep a stool diary if changes are intermittent
\end{itemize}

\section*{How Your Doctor Will Evaluate You}
\begin{itemize}
\item History and physical examination (including digital rectal exam)
\item Fecal occult blood test
\item Complete blood count (check for anemia)
\item Liver function tests and bilirubin
\item Colonoscopy for suspected lower GI source
\item Upper endoscopy for suspected upper GI bleeding
\end{itemize}

\section*{Treatment Options}
\begin{itemize}
\item Treat underlying cause
\item For GI bleeding: endoscopic therapy, IV fluids, transfusions
\item For biliary blockage: ERCP with stenting, surgery
\item For malabsorption: dietary changes, enzyme replacement
\item Discontinue or adjust causative medications
\end{itemize}
